\section{Code--Data Identification and Data-Preserving Rewriting}
\label{sec:code_data}
\subsection{Why Data Identification Comes Before Transformation}
A raw binary does not tell a rewriting tool which bytes are instructions
and which are data. A short program may contain both executable code and
an embedded string or lookup table. A disassembler translates bytes into
assembly instructions, but the ability to decode a byte sequence does not
establish that the processor will execute it. String bytes may also decode
as plausible instructions. Rewriting those apparent instructions can
therefore corrupt a value that the program subsequently reads.

This is the failure illustrated by the inline string in
Table~\ref{tab:broken-payload}. It differs from
Table~\ref{tab:functional-payload}, where the string is encoded inside an
instruction operand: there, the surrounding instruction must remain
intact. A useful identification method must distinguish these situations,
rather than marking every printable sequence as a separate data object.

The extension introduces a preparation stage with two responsibilities.
\emph{Identification} proposes where data occurs. \emph{Data handling}
decides which proposed objects the engine can retain through a layout
change, preserves their contents, and updates supported references
to their new locations. A reference is an instruction operand or address
calculation that points to a code or data location. Preserving a string's
bytes is insufficient if the program subsequently reads its old address.

Figure~\ref{fig:extended_pipeline} locates these responsibilities in the
extended engine. The symbolic/heuristic and LLM classifiers are alternative
ways to propose data regions. Separately, the transformation itself can be
implemented by the native engine or proposed by an LLM. Thus, using an LLM
to identify data does not imply using it to rewrite instructions.

\begin{figure}[t]
\centering
\begin{tikzpicture}[
  box/.style={draw,rounded corners=2pt,align=center,font=\small,
    minimum height=0.8cm,text width=4.0cm},
  arrow/.style={-{Latex[length=2mm]},thick}]
\node[box,text width=5.7cm] (input) at (0,0) {Raw binary and target architecture};
\node[box] (sym) at (-2.35,-1.3) {Symbolic exploration\\and static heuristics};
\node[box] (llm) at (2.35,-1.3) {LLM data classifier\\(proposes regions only)};
\node[box,text width=7.4cm] (data) at (0,-2.65) {Shared data handling: check candidate regions,\\preserve supported data and track references};
\node[box,text width=7.4cm] (trans) at (0,-4.0) {Transform code: native methods, assembly wrapping,\\or LLM-proposed EQ, NOP, BCF, or IBP};
\node[box,text width=7.4cm] (build) at (0,-5.3) {Finalize layout and references; check construction};
\node[box,text width=7.4cm] (test) at (0,-6.6) {Test the produced binary with an external harness};
\draw[arrow] (input) -- (sym);
\draw[arrow] (input) -- (llm);
\draw[arrow] (sym) -- (data);
\draw[arrow] (llm) -- (data);
\draw[arrow] (data) -- (trans);
\draw[arrow] (trans) -- (build);
\draw[arrow] (build) -- (test);
\node[font=\small,fill=white] at (0,-1.3) {or};
\end{tikzpicture}
\caption{Conceptual organization of the extensions. Identification chooses
candidate data regions; shared handling preserves supported objects during
rewriting. Successful construction is followed by a separate behavioral
test. Unsupported inputs can stop before a binary is produced.}
\label{fig:extended_pipeline}
\end{figure}

\subsection{What a Classification Result Means}
For an input of $n$ bytes, a proposed data region is an interval $[s,e)$:
it starts at offset $s$ from the beginning of the input and includes every
byte up to, but not including, $e$. For example, $[2,8)$ denotes six bytes.
An optional type describes a string encoding, such as ASCII or UTF-16LE,
or another data object. The output is a set of candidate intervals,
not a replacement binary.

Three stages must be distinguished: a classifier's proposed regions,
regions accepted by the engine's checks, and regions actually extracted
from the instruction stream. A correct proposal may still be unsupported
if its references cannot be updated. Conversely, different proposals can
produce the same binary when the shared checks retain the same objects.
The experiments compare annotations, accepted outputs, and final hashes
at their respective stages.

\subsection{Symbolic Exploration Combined with Static Heuristics}
The non-LLM method combines three sources of evidence. First, static
analysis recovers a control-flow graph (CFG): nodes represent basic blocks
of instructions and edges represent possible transfers of execution.
Zone uses angr's CFGFast analysis for this initial structure~\cite{angr_cfg_docs}.
Second, symbolic execution explores possible paths while representing
unknown register and memory values as symbolic expressions. This supplies
additional evidence about which instructions can be reached. Third,
pattern-based rules look for candidate strings and graph-inferred data.

An observed instruction protects its entire decoded byte span, not only
its starting address. This matters for variable-length instructions: an
interior byte should not become a data candidate merely because the
explorer recorded only the instruction's first byte. Newly observed entry
points can also support another graph-recovery pass. Printable candidates
that overlap observed code are excluded by the string heuristic; candidates
overlapping graph-recovered code need additional structural support.

Exploration is bounded by time, solver, and state limits. The initial state
does not reproduce every feature of a Windows process, including all
external API behavior. Consequently, an unvisited region is \emph{unknown},
not proven to be data. Reports retain the evidence category, stopping
reason, unresolved states, and errors. These distinctions prevent a
partial execution search from being presented as a complete partition of
the binary into code and data.

\subsection{LLM Identification and Shared Data Handling}
The alternative classifier receives the raw bytes in hexadecimal, a
linear disassembly, and candidate printable strings. Linear disassembly
means decoding instructions in byte order; its listing is evidence for the
model to assess, not a trusted annotation. The requested answer contains
structured data intervals. The model does not perform extraction or
rewrite the program in this mode.

Deterministic parsing rejects intervals outside the input, empty intervals,
and invalid overlaps. The engine then examines supported references and
applies its data-handling rules to a fresh working instance. Accepted data
contents are retained separately from the instruction representation;
supported references are finalized after transformation, when the new
locations are known. Integrity checks verify the preserved contents.
Architecture-specific addressing requires separate x64 and x86 handling.

For a benign example, consider a routine that copies the embedded strings
\texttt{alpha} and \texttt{beta} into an output buffer. Identification finds
the two string regions; handling preserves both strings and their links
to the copying code; transformation changes the instruction representation.
The final test checks that the intended strings reach the intended buffer
positions. Correct interval identification, successful reference handling,
and correct observed output are three distinct questions in this example.

The general classifier can receive validation feedback and make another
proposal. Its worked-example setting is independent of the zero-shot and
two-shot transformation conditions defined next. Cached identification can
reuse regions associated with an input identity, while references are
recomputed for the working representation. These shared steps also explain
why agreement between final binaries is not, by itself, a measure of the
LLM classifier's accuracy.
